% chapter3_verification.tex
% 第3章 3.2节 检验有效性的蒙特卡洛验证设计

\section{检验有效性的蒙特卡洛验证设计}\label{sec:mc-design}

在高维 VAR 模型中引入 Lasso 稀疏惩罚或低秩约束后，似然比统计量在有限样本下的分布偏离经典渐近理论，无法通过解析方法获取精确临界值。为验证上节所构建的残差 Bootstrap 检验方案在有限样本下的统计特性，本节设计了一套系统的蒙特卡洛仿真实验。其核心逻辑是：在已知数据生成机制（Data Generating Process, DGP）的受控环境中，通过大规模重复试验，对检验方案的误差控制能力与信号识别能力进行量化评估。

整个验证体系围绕两个核心指标展开：第一类错误率（Size）和统计功效（Power）。前者考察在原假设成立时，检验是否将错误拒绝的概率控制在名义显著性水平附近；后者考察在备择假设成立时，检验能否有效识别真实的结构性变化。

为确保结论的普适性，验证设计建立了由低维至高维、由无约束至正则化约束逐层递进的三层实验架构：
\begin{itemize}
    \item \textbf{第一层：低维基准层}（$N=10$，OLS 估计）。在 OLS 可行的标准场景中，将 Bootstrap 似然比检验与经典渐近 $F$ 检验进行对比，验证二者在 Size 控制与功效表现上的一致性，为重抽样机制的正确性建立信任锚点。
    \item \textbf{第二层：高维稀疏约束层}（$N=20$，Lasso 估计）。在系数矩阵仅有少数非零元素的稀疏场景中，检验 Lasso 配合 Bootstrap 似然比推断是否能提供有效的断裂检测路径。
    \item \textbf{第三层：高维低秩约束层}（$N=20$，降秩回归估计）。在系数矩阵由少数潜在因子驱动的低秩场景中，检验降秩回归（Reduced-Rank Regression, RRR）配合 Bootstrap 似然比推断的有效性。
\end{itemize}

三层实验共享相同的假设定义（$H_0: \Phi_1 = \Phi_2$）和样本量口径，差异仅在于参数估计器与 $p$ 值计算方式。表~\ref{tab:mc-models} 列出了四类模型的配置参数。

\begin{table}[htbp]
\centering
\caption{蒙特卡洛仿真实验的四类模型配置}\label{tab:mc-models}
\begin{tabular}{lcccccc}
\toprule
模型 & 估计器 & $N$ & DGP 类型 & $p$ 值方法 & 扰动方向 \\
\midrule
baseline\_ols\_f & OLS & 10 & 稠密 & 渐近 $F$ 检验 & 均匀 \\
baseline\_ols    & OLS & 10 & 稠密 & Bootstrap LR & 均匀 \\
sparse\_lasso    & Lasso & 20 & 稀疏（$s=0.15$） & Bootstrap LR & 稀疏支撑集 \\
lowrank\_rrr     & RRR & 20 & 低秩（$r=2$） & Bootstrap LR & 低秩子空间 \\
\bottomrule
\end{tabular}
\end{table}

所有实验的公共参数设定如下：样本总长度 $T = 500$，滞后阶数 $p = 1$，已知断点位于序列中点 $t^* = 250$，残差协方差矩阵 $\Sigma = 0.5 \cdot I_N$，Bootstrap 重抽样次数 $B = 500$，名义显著性水平 $\alpha = 0.05$。每组实验在两个独立随机种子下分别运行，最终结果取双种子平均值以降低随机波动。

\subsection{第一类错误率的蒙特卡洛评估方案}\label{subsec:type1}

第一类错误率衡量的是当原假设完全成立时，检验错误拒绝 $H_0$ 的概率。在蒙特卡洛仿真中，这一指标被量化为经验拒绝率（Empirical Rejection Rate）。对于高维模型，Lasso 的变量选择偏差或 RRR 的低秩投影偏差会在对数似然的差分运算中被累积放大，导致似然比统计量的分布右尾增厚。若 Bootstrap 经验分布无法准确捕捉这一偏移，经验拒绝率将显著超过名义水平 $\alpha$，产生过度拒绝现象。因此，Size 评估的核心目的在于验证残差 Bootstrap 经验分布能否在有限样本下有效校正高维估计偏差，将第一类错误率稳定控制在 $\alpha$ 附近。

\subsubsection{数据生成过程}

在 $H_0$ 下，整条时间序列由单一系数矩阵 $\Phi_1$ 驱动，不包含任何结构性变化。系数矩阵根据实验层级采用不同的生成策略：

\paragraph{低维稠密（第一层，$N=10$）.}
矩阵元素独立采样自 $\Phi_{ij} \sim \mathcal{N}(0, \mathrm{scale}^2)$，其中 $\mathrm{scale} = \min(0.3, \; 0.85/\sqrt{N})$。反复生成直至满足平稳性条件~\eqref{eq:stationarity}。

\paragraph{高维稀疏（第二层，$N=20$）.}
首先按上述方式生成稠密矩阵，随后以概率 $1 - s$（$s = 0.15$）将元素置零，使约 $85\%$ 的元素为零，模拟经济网络中仅有少数核心节点具备显著传导效应的稀疏连接结构。

\paragraph{高维低秩（第三层，$N=20$）.}
通过低秩分解 $\Phi_1 = UV^\top$ 生成，其中 $U \in \mathbb{R}^{N \times r}$，$V \in \mathbb{R}^{Np \times r}$，$r = 2$，各元素独立采样自 $\mathcal{N}(0, 0.3^2)$。生成后将 $\Phi_1$ 缩放至目标谱半径 $\rho(\Phi_1) = 0.40$，以消除不同随机种子间因随机生成导致的谱半径方差。

残差协方差矩阵固定为 $\Sigma = 0.5 \cdot I_N$，设定为对角阵以剥离误差项截面相关性对检验结果的混淆影响，使评估焦点集中于系数矩阵变动的检测。

\subsubsection{仿真流程}

在每次蒙特卡洛迭代中，算法首先通过多元正态分布生成独立同分布白噪声序列，结合 $\Phi_1$ 按 VAR($p$) 递推方程逐期生成长度为 $T + T_{\text{burn}}$ 的时间序列，其中 $T_{\text{burn}} = 100$ 为预热期（Burn-in period），用于消除随机初始值对平稳分布的影响。截去预热期后，保留 $T$ 期有效观测。

对该无断裂样本，按算法~\ref{alg:bootstrap-test} 执行完整的 Bootstrap 似然比检验流程（对于渐近 $F$ 模型，则按经典 Chow 检验的 $F$ 分布计算 $p$ 值），记录是否拒绝 $H_0$。

为考察经验拒绝率随蒙特卡洛重复次数的收敛行为，Size 评估在网格 $M \in \{50, 100, 300, 500, 1000, 2000\}$ 上分别计算经验第一类错误率：
\begin{equation}\label{eq:type1-rate}
    \hat{\alpha}(M) = \frac{1}{M} \sum_{m=1}^{M} \mathbf{1}\{\hat{p}_m \leq \alpha\},
\end{equation}
其中 $\hat{p}_m$ 为第 $m$ 次迭代的经验 $p$ 值。$\hat{\alpha}(M)$ 与名义水平 $\alpha$ 之差即为尺寸失真（Size distortion）的度量。若 $\hat{\alpha}(M)$ 随 $M$ 增大收敛至 $\alpha = 0.05$ 附近的窄区间，则表明 Bootstrap 经验分布有效校正了高维估计偏差，检验具有良好的 Size 控制能力。算法~\ref{alg:type1-mc} 给出了第一类错误率蒙特卡洛评估的完整流程。

\begin{algorithm}[htbp]
\caption{第一类错误率的蒙特卡洛评估}\label{alg:type1-mc}
\KwIn{系数矩阵 $\Phi_1$，协方差 $\Sigma$，样本量 $T$，断点 $t^*$，滞后阶 $p$，蒙特卡洛次数 $M$，Bootstrap 次数 $B$，显著性水平 $\alpha$}
\KwOut{经验第一类错误率 $\hat{\alpha}$}
\BlankLine
$\mathrm{reject\_count} \leftarrow 0$\;
\For{$m = 1, 2, \ldots, M$}{
    \tcc{数据生成（$H_0$ 下，无断点）}
    生成白噪声 $\{\varepsilon_t\}_{t=1}^{T+T_{\text{burn}}} \sim \mathcal{N}(0, \Sigma)$\;
    以 $\Phi_1$ 为系数矩阵，按 VAR($p$) 递推生成序列，截去前 $T_{\text{burn}}$ 期\;
    \BlankLine
    \tcc{检验执行}
    对全样本按对应估计器拟合约束模型，计算 $\ell_0$\;
    以 $t^*$ 为断点分段拟合非约束模型，计算 $\ell_1$\;
    $\mathrm{LR}_{\text{obs}} \leftarrow 2(\ell_1 - \ell_0)$\;
    \BlankLine
    \tcc{$p$ 值计算}
    \eIf{使用 Bootstrap}{
        按算法~\ref{alg:bootstrap-test} 步骤 2--4 计算经验 $p$ 值 $\hat{p}$\;
    }{
        按渐近 $F$ 分布计算 $p$ 值 $\hat{p}$\;
    }
    \If{$\hat{p} \leq \alpha$}{
        $\mathrm{reject\_count} \leftarrow \mathrm{reject\_count} + 1$\;
    }
}
$\hat{\alpha} \leftarrow \mathrm{reject\_count} \;/\; M$\;
\end{algorithm}


\subsection{统计功效的蒙特卡洛评估方案}\label{subsec:power}

当 Size 评估确认检验方案在 $H_0$ 下具有良好的错误控制能力后，需进一步考察其在 $H_1$ 下的信号识别能力。统计功效衡量的是当系数矩阵确实在断点处发生变化时（即 $\Phi_1 \neq \Phi_2$），检验正确拒绝 $H_0$ 的概率。在高维系统中，真实的结构变化往往表现为少数网络连接权重的漂移或潜在因子载荷的旋转，这些局部信号容易被高维噪声淹没，或被正则化估计器作为冗余信息过滤。因此，功效评估需通过逐步增大结构偏离幅度，考察检验识别能力是否随信号强度单调递增。

\subsubsection{效应量的定义与量化}

为在不同维度与不同结构约束的实验层级之间建立统一且可比的信号强度度量，本研究引入效应量（Effect Size）$\delta$，定义为断裂前后系数矩阵差的 Frobenius 范数：
\begin{equation}\label{eq:effect-size}
    \delta = \|\Phi_2 - \Phi_1\|_F = \left(\sum_{i,j} (\Phi_{2,ij} - \Phi_{1,ij})^2\right)^{1/2}.
\end{equation}
Frobenius 范数将高维矩阵空间中的参数位移压缩为单一非负标量，等价于差异矩阵全部奇异值平方和的平方根，直观地度量了结构冲击所释放的全局信号能量。通过逐步增大 $\delta$ 的值，可以构造从微弱扰动到剧烈断裂的一系列备择假设场景。本实验采用 $\delta \in \{0.1, 0.3, 0.5, 1.0\}$ 的效应量网格，覆盖从接近 Size 水平到功效趋近于 $1$ 的完整检测区间。

\subsubsection{双区制数据生成机制}

在备择假设下，时间序列在断点 $t^*$ 处被切分为两个区制。断点前 $t < t^*$ 的数据由 $\Phi_1$ 驱动，断点后 $t \geq t^*$ 的数据由 $\Phi_2$ 接管：
\begin{equation}\label{eq:two-regime}
    Y_t = \begin{cases}
        c + \Phi_1 \cdot \mathrm{vec}(Y_{t-1}, \ldots, Y_{t-p}) + \varepsilon_t, & t < t^*, \\
        c + \Phi_2 \cdot \mathrm{vec}(Y_{t-1}, \ldots, Y_{t-p}) + \varepsilon_t, & t \geq t^*.
    \end{cases}
\end{equation}
后半段的 VAR 递推以前半段末尾的 $p$ 个观测值作为初始条件，保证样本路径的连续性。断裂后的系数矩阵构造为
\begin{equation}\label{eq:phi2-construct}
    \Phi_2 = \Phi_1 + \delta \cdot D,
\end{equation}
其中 $D$ 为经 Frobenius 范数归一化的扰动方向矩阵（$\|D\|_F = 1$），其具体形式取决于实验层级。

\subsubsection{结构匹配扰动方向}\label{subsubsec:perturbation}

在功效评估中，扰动方向 $D$ 的设计至关重要。若对稀疏系统施加稠密扰动，$\Phi_2$ 将丧失稀疏性，Lasso 估计器面临的将不再是参数值的偏移，而是模型结构本身的崩溃；类似地，若对低秩系统施加满秩扰动，$\Phi_2$ 的秩将膨胀至满秩，降秩回归估计器无法在有效子空间中工作。为保证功效评估的纯粹性——即检验识别的是参数数值的偏移而非结构类型的改变——本研究采用与各层级数据生成机制相匹配的扰动方向：

\paragraph{第一层（低维稠密）.}
扰动方向为均匀全 $1$ 矩阵，经范数归一化：
\begin{equation}
    D = \frac{\mathbf{1}_{N \times Np}}{\|\mathbf{1}_{N \times Np}\|_F}.
\end{equation}
该设计不偏向任何特定参数位置，考察全域敏感性。

\paragraph{第二层（高维稀疏）.}
扰动仅作用于 $\Phi_1$ 的非零支撑集上，通过二元掩码矩阵实现：
\begin{equation}
    D = \frac{\mathrm{support}(\Phi_1)}{\|\mathrm{support}(\Phi_1)\|_F}, \quad \mathrm{support}(\Phi_1)_{ij} = \mathbf{1}\{|\Phi_{1,ij}| > 10^{-10}\}.
\end{equation}
该掩码将所有试图在零元素位置建立新连接的扰动信号切除，确保断裂前后的稀疏结构严格一致。断裂仅表现为已有传导路径上信号强度的变化，而非新路径的产生。

\paragraph{第三层（高维低秩）.}
扰动方向被约束在 $\Phi_1$ 的前 $r$ 个奇异向量张成的子空间内。对 $\Phi_1$ 进行奇异值分解 $\Phi_1 = U \Sigma_s V^\top$，取前 $r$ 个左、右奇异向量 $U_r$、$V_r$，构造
\begin{equation}
    D = \frac{U_r \mathbf{1}_{r \times r} V_r^\top}{\|U_r \mathbf{1}_{r \times r} V_r^\top\|_F}.
\end{equation}
该投影保证 $\Phi_2$ 的秩不超过 $r$，断裂仅表现为低维因子空间内载荷大小的变化，而非因子维度的膨胀。

上述结构匹配设计基于实证中的观察：真实世界的结构断裂通常是保结构的——稀疏系统的断裂改变的是已有连接的强度，低秩系统的断裂改变的是因子载荷的大小。

\subsubsection{平稳性校验与收缩机制}

在将扰动 $\delta \cdot D$ 叠加至 $\Phi_1$ 后，需验证 $\Phi_2$ 是否满足平稳性条件~\eqref{eq:stationarity}。即便扰动受到结构匹配约束，在高维矩阵空间中，线性叠加仍可能将伴随矩阵的谱半径推出单位圆。若 $\Phi_2$ 不满足平稳性，由其驱动的后半段序列将呈指数发散，对数似然函数与残差同方差假设均不再成立，功效评估失去意义。

为此，本研究在 $\Phi_2$ 的生成过程中嵌入闭环校验机制：每次构造候选 $\Phi_2$ 后，立即计算其伴随矩阵的最大特征值模。若谱半径超过 $1$，则将扰动乘以收缩因子 $\gamma = 0.9$ 进行等比压缩，重新叠加并再次校验：
\begin{equation}\label{eq:shrinkage}
    \Phi_2^{(k)} = \Phi_1 + \delta \cdot \gamma^k \cdot D, \quad k = 1, 2, \ldots, K_{\max},
\end{equation}
其中 $K_{\max} = 30$。该反馈循环持续执行直至获得一个既保持目标结构又严格平稳的合法系数矩阵。由于收缩操作可能使实际 Frobenius 距离小于名义效应量 $\delta$，实验同时记录名义值与实际值以供分析。

\subsubsection{功效计算}

固定蒙特卡洛重复次数 $M = 500$，对每个效应量 $\delta$，功效定义为
\begin{equation}\label{eq:power}
    \widehat{\mathrm{power}}(\delta) = \frac{1}{M} \sum_{m=1}^{M} \mathbf{1}\{\hat{p}_m \leq \alpha\},
\end{equation}
其中 $\hat{p}_m$ 为第 $m$ 条含断裂序列的经验 $p$ 值。若功效随 $\delta$ 的增大呈单调递增趋势，且在 $\delta$ 足够大时趋近于 $1$，则表明检验方案具备有效的判别能力。算法~\ref{alg:power-mc} 给出了统计功效蒙特卡洛评估的完整流程。

\begin{algorithm}[htbp]
\caption{统计功效的蒙特卡洛评估}\label{alg:power-mc}
\KwIn{系数矩阵 $\Phi_1$，协方差 $\Sigma$，效应量网格 $\{\delta_1, \ldots, \delta_K\}$，样本量 $T$，断点 $t^*$，滞后阶 $p$，蒙特卡洛次数 $M$，Bootstrap 次数 $B$，显著性水平 $\alpha$}
\KwOut{功效曲线 $\{\widehat{\mathrm{power}}(\delta_k)\}_{k=1}^K$}
\BlankLine
\For{$k = 1, 2, \ldots, K$}{
    \tcc{构造断裂后系数矩阵}
    根据实验层级生成结构匹配的扰动方向 $D$（$\|D\|_F = 1$）\;
    $\Phi_2 \leftarrow \Phi_1 + \delta_k \cdot D$\;
    \While{$\Phi_2$ 不满足平稳性条件}{
        $D \leftarrow 0.9 \cdot D$，$\Phi_2 \leftarrow \Phi_1 + \delta_k \cdot D$\;
    }
    \BlankLine
    $\mathrm{reject\_count} \leftarrow 0$\;
    \For{$m = 1, 2, \ldots, M$}{
        \tcc{双区制数据生成（$H_1$ 下，含断点）}
        以 $\Phi_1$ 驱动生成 $t < t^*$ 的前半段序列（含预热期，事后截去）\;
        以 $\Phi_2$ 驱动生成 $t \geq t^*$ 的后半段序列\;
        \BlankLine
        \tcc{检验执行与 $p$ 值计算}
        按算法~\ref{alg:bootstrap-test} 执行完整检验，得到经验 $p$ 值 $\hat{p}$\;
        \If{$\hat{p} \leq \alpha$}{
            $\mathrm{reject\_count} \leftarrow \mathrm{reject\_count} + 1$\;
        }
    }
    $\widehat{\mathrm{power}}(\delta_k) \leftarrow \mathrm{reject\_count} \;/\; M$\;
}
\end{algorithm}

通过上述设计，蒙特卡洛验证方案构建了一个不依赖解析分布的数值实验平台，能够在高维正则化约束等非标准条件下定量评估 Bootstrap 似然比检验的水平控制能力和检验功效，为后续实证应用中判断检验工具的适用范围提供依据。
